% Q12_answers.tex — your answers to paper/Q12_problems.tex (learner-owned). One section
% per problem, keyed by its \label{ans:B<n>} anchor; the GRADER comment convention and the
% PREREG markers are paper/template_answers.tex's. /grade-paper hashes exactly the lines
% from PREREG-BEGIN to PREREG-END, inclusive: editing anything between them after grading
% changes the hash, and the runs stay locked.
\documentclass[11pt]{article}
% preamble.tex — shared preamble for the dafny-verus-learn-0 written track.
%
% Usage (every <ID>_problems.tex, <ID>_answers.tex and solutions/<ID>_solutions.tex):
%
%   \documentclass[11pt]{article}
%   % preamble.tex — shared preamble for the dafny-verus-learn-0 written track.
%
% Usage (every <ID>_problems.tex, <ID>_answers.tex and solutions/<ID>_solutions.tex):
%
%   \documentclass[11pt]{article}
%   % preamble.tex — shared preamble for the dafny-verus-learn-0 written track.
%
% Usage (every <ID>_problems.tex, <ID>_answers.tex and solutions/<ID>_solutions.tex):
%
%   \documentclass[11pt]{article}
%   \input{preamble}          % the Makefile always runs latexmk from paper/
%   \begin{document}
%   ...
%   \end{document}
%
% Engine: pdflatex. Dafny and Verus sources are mostly ASCII; the `literate`
% tables below map the common Unicode symbols to math glyphs inside listings.
% For other Unicode, extend a table, or build that file with lualatex
% (`latexmk -lualatex <file>`) after replacing the inputenc/fontenc lines with
% \usepackage{fontspec} and a monospace font with good coverage.
%
% Contents: program-logic macros (Hoare triples, wp, VCs, substitution,
% semantics brackets, typing judgments), inference rules (mathpartir),
% theorem and problem environments, the answer environment, the grader
% margin-note macro, and listings languages `dafny` and `verus` (Python uses
% listings' built-in `Python`), with inline forms \dfy{...}, \vrs{...}, \py{...}.
\usepackage[utf8]{inputenc}
\usepackage[T1]{fontenc}

\usepackage{amsmath}
\usepackage{amssymb}
\usepackage{amsthm}
\usepackage{stmaryrd}         % \llbracket \rrbracket for semantics brackets
\usepackage{mathpartir}       % inference rules: \inferrule, the mathpar environment
\usepackage{listings}
\usepackage{xcolor}
\usepackage[margin=1.1in]{geometry}

% ---------------------------------------------------------------------------
% Program logic.
%   \hoare{P}{c}{Q}       {P} c {Q}           (partial correctness)
%   \thoare{P}{c}{Q}      [P] c [Q]           (total correctness)
%   \provable{P}{c}{Q}    |- {P} c {Q}        (derivable in the proof system)
%   \WP{c}{Q}             wp(c, Q)            (\wp is LaTeX's Weierstrass p, so
%   \WLP{c}{Q}            wlp(c, Q)            the macro names are upper case)
%   \VC{c}{Q}             VC(c, Q)
%   \subst{Q}{x}{e}       Q[e/x]              (e replaces x in Q)
%   \sem{e}               [[e]]               (denotation)
%   \assign, \seqc        :=  and the sequencing semicolon
%   \impl                 ==> as a math arrow (LaTeX's own \iff gives <==>)
% ---------------------------------------------------------------------------
\newcommand{\hoare}[3]{\{\,#1\,\}\;#2\;\{\,#3\,\}}
\newcommand{\thoare}[3]{[\,#1\,]\;#2\;[\,#3\,]}
\newcommand{\provable}[3]{\vdash \hoare{#1}{#2}{#3}}
\newcommand{\WP}[2]{\mathrm{wp}(#1,\,#2)}
\newcommand{\WLP}[2]{\mathrm{wlp}(#1,\,#2)}
\newcommand{\VC}[2]{\mathrm{VC}(#1,\,#2)}
\newcommand{\subst}[3]{#1[#3/#2]}
\newcommand{\sem}[1]{\llbracket #1 \rrbracket}
\newcommand{\assign}{\mathrel{:=}}
\newcommand{\seqc}{\mathbin{;}\,}
\newcommand{\impl}{\Longrightarrow}

% ---------------------------------------------------------------------------
% Judgments.  \judgment{\Gamma}{e}{\tau}  ==>  Gamma |- e : tau
% ---------------------------------------------------------------------------
\newcommand{\judgment}[3]{#1 \vdash #2 : #3}
\newcommand{\judg}[3]{\judgment{#1}{#2}{#3}}          % short alias
\newcommand{\emptyjudgment}[2]{\cdot \vdash #1 : #2}  % empty context
\newcommand{\ctx}{\Gamma}
\newcommand{\turnstile}{\vdash}

% ---------------------------------------------------------------------------
% Theorem environments.
% ---------------------------------------------------------------------------
\theoremstyle{plain}
\newtheorem{theorem}{Theorem}[section]
\newtheorem{lemma}[theorem]{Lemma}
\newtheorem{corollary}[theorem]{Corollary}
\newtheorem{proposition}[theorem]{Proposition}
\theoremstyle{definition}
\newtheorem{definition}[theorem]{Definition}
% Written-track problems are numbered B1, B2, ... on both tracks (P<n> is a
% predict item). Label each one \label{prob:B<n>} and its answer \label{ans:B<n>}.
\newtheorem{problem}{Problem}
\renewcommand{\theproblem}{B\arabic{problem}}
\theoremstyle{remark}
\newtheorem*{remark}{Remark}
\newenvironment{answer}
  {\par\medskip\noindent\textbf{Answer.}\begin{quotation}\noindent}
  {\end{quotation}\medskip}

% ---------------------------------------------------------------------------
% Grader margin notes. Rendered form of the `% GRADER[...][...][...]:` comment
% convention (see template_answers.tex). Use inline where a visible note helps:
%   \gradernote{B2/induction-soundness/2: the invariant is not preserved by the last assignment}
% ---------------------------------------------------------------------------
\definecolor{gradercolor}{RGB}{160, 30, 30}
\newcommand{\gradernote}[1]{\textcolor{gradercolor}{\small\sffamily [GRADER: #1]}}

% ---------------------------------------------------------------------------
% Code listings: Dafny, Verus (Rust inside verus!{...}), Python.
% ---------------------------------------------------------------------------
\definecolor{codekeyword}{RGB}{0, 60, 130}
\definecolor{codecomment}{RGB}{90, 110, 90}
\definecolor{codestring}{RGB}{130, 60, 10}
\definecolor{codehole}{RGB}{180, 0, 0}

\lstdefinelanguage{dafny}{
  sensitive=true,
  morekeywords={
    abstract, allocated, as, assert, assume, bool, break, by, calc, case, char,
    class, codatatype, const, constructor, continue, datatype, decreases, else,
    ensures, exists, expect, export, extends, false, forall, fresh, function,
    ghost, if, import, in, include, int, invariant, is, iterator, label, lemma,
    map, match, method, modifies, modify, module, multiset, nat, new, newtype,
    null, old, opaque, opened, predicate, print, provides, reads, real, refines,
    requires, return, returns, reveal, reveals, seq, set, static, string, then,
    this, trait, true, twostate, type, unchanged, var, while, witness},
  emph={assume},
  emphstyle=\color{codehole}\bfseries,
  comment=[l]{//},
  morecomment=[s]{/*}{*/},
  morestring=[b]",
}

\lstdefinelanguage{verus}{
  sensitive=true,
  morekeywords={
    as, break, const, continue, crate, dyn, else, enum, false, fn, for, if,
    impl, in, let, loop, match, mod, move, mut, pub, ref, return, self, Self,
    static, struct, super, trait, true, type, unsafe, use, where, while,
    spec, proof, exec, open, closed, tracked, ghost, requires, ensures,
    recommends, invariant, decreases, assert, forall, exists, choose, by,
    broadcast, uninterp, axiom, returns, int, nat, Seq, Set, Map, Ghost, Tracked},
  emph={assume, admit},
  emphstyle=\color{codehole}\bfseries,
  comment=[l]{//},
  morecomment=[s]{/*}{*/},
  morestring=[b]",
}

\lstset{
  basicstyle=\ttfamily\small,
  keywordstyle=\color{codekeyword}\bfseries,
  commentstyle=\color{codecomment}\itshape,
  stringstyle=\color{codestring},
  showstringspaces=false,
  columns=fullflexible,
  keepspaces=true,
  breaklines=true,
  frame=leftline,
  framerule=0.6pt,
  xleftmargin=1em,
  aboveskip=0.8em,
  belowskip=0.8em,
  literate=
    {→}{{$\to$}}1        {←}{{$\leftarrow$}}1  {⇒}{{$\Rightarrow$}}1
    {⇔}{{$\Leftrightarrow$}}1                  {↦}{{$\mapsto$}}1
    {∀}{{$\forall$}}1    {∃}{{$\exists$}}1     {λ}{{$\lambda$}}1
    {∧}{{$\wedge$}}1     {∨}{{$\vee$}}1        {¬}{{$\neg$}}1
    {≠}{{$\neq$}}1       {≤}{{$\leq$}}1        {≥}{{$\geq$}}1
    {∈}{{$\in$}}1        {∉}{{$\notin$}}1      {⊆}{{$\subseteq$}}1
    {×}{{$\times$}}1     {·}{{$\cdot$}}1       {∅}{{$\varnothing$}}1
    {ℕ}{{$\mathbb{N}$}}1 {ℤ}{{$\mathbb{Z}$}}1
    {⊢}{{$\vdash$}}1     {⟦}{{$\llbracket$}}1  {⟧}{{$\rrbracket$}}1,
}

% Code environments and inline code.
\lstnewenvironment{dafnycode}{\lstset{language=dafny}}{}
\lstnewenvironment{veruscode}{\lstset{language=verus}}{}
\lstnewenvironment{pycode}{\lstset{language=Python}}{}
\newcommand{\dfy}[1]{\lstinline[language=dafny]{#1}}
\newcommand{\vrs}[1]{\lstinline[language=verus]{#1}}
\newcommand{\py}[1]{\lstinline[language=Python]{#1}}
          % the Makefile always runs latexmk from paper/
%   \begin{document}
%   ...
%   \end{document}
%
% Engine: pdflatex. Dafny and Verus sources are mostly ASCII; the `literate`
% tables below map the common Unicode symbols to math glyphs inside listings.
% For other Unicode, extend a table, or build that file with lualatex
% (`latexmk -lualatex <file>`) after replacing the inputenc/fontenc lines with
% \usepackage{fontspec} and a monospace font with good coverage.
%
% Contents: program-logic macros (Hoare triples, wp, VCs, substitution,
% semantics brackets, typing judgments), inference rules (mathpartir),
% theorem and problem environments, the answer environment, the grader
% margin-note macro, and listings languages `dafny` and `verus` (Python uses
% listings' built-in `Python`), with inline forms \dfy{...}, \vrs{...}, \py{...}.
\usepackage[utf8]{inputenc}
\usepackage[T1]{fontenc}

\usepackage{amsmath}
\usepackage{amssymb}
\usepackage{amsthm}
\usepackage{stmaryrd}         % \llbracket \rrbracket for semantics brackets
\usepackage{mathpartir}       % inference rules: \inferrule, the mathpar environment
\usepackage{listings}
\usepackage{xcolor}
\usepackage[margin=1.1in]{geometry}

% ---------------------------------------------------------------------------
% Program logic.
%   \hoare{P}{c}{Q}       {P} c {Q}           (partial correctness)
%   \thoare{P}{c}{Q}      [P] c [Q]           (total correctness)
%   \provable{P}{c}{Q}    |- {P} c {Q}        (derivable in the proof system)
%   \WP{c}{Q}             wp(c, Q)            (\wp is LaTeX's Weierstrass p, so
%   \WLP{c}{Q}            wlp(c, Q)            the macro names are upper case)
%   \VC{c}{Q}             VC(c, Q)
%   \subst{Q}{x}{e}       Q[e/x]              (e replaces x in Q)
%   \sem{e}               [[e]]               (denotation)
%   \assign, \seqc        :=  and the sequencing semicolon
%   \impl                 ==> as a math arrow (LaTeX's own \iff gives <==>)
% ---------------------------------------------------------------------------
\newcommand{\hoare}[3]{\{\,#1\,\}\;#2\;\{\,#3\,\}}
\newcommand{\thoare}[3]{[\,#1\,]\;#2\;[\,#3\,]}
\newcommand{\provable}[3]{\vdash \hoare{#1}{#2}{#3}}
\newcommand{\WP}[2]{\mathrm{wp}(#1,\,#2)}
\newcommand{\WLP}[2]{\mathrm{wlp}(#1,\,#2)}
\newcommand{\VC}[2]{\mathrm{VC}(#1,\,#2)}
\newcommand{\subst}[3]{#1[#3/#2]}
\newcommand{\sem}[1]{\llbracket #1 \rrbracket}
\newcommand{\assign}{\mathrel{:=}}
\newcommand{\seqc}{\mathbin{;}\,}
\newcommand{\impl}{\Longrightarrow}

% ---------------------------------------------------------------------------
% Judgments.  \judgment{\Gamma}{e}{\tau}  ==>  Gamma |- e : tau
% ---------------------------------------------------------------------------
\newcommand{\judgment}[3]{#1 \vdash #2 : #3}
\newcommand{\judg}[3]{\judgment{#1}{#2}{#3}}          % short alias
\newcommand{\emptyjudgment}[2]{\cdot \vdash #1 : #2}  % empty context
\newcommand{\ctx}{\Gamma}
\newcommand{\turnstile}{\vdash}

% ---------------------------------------------------------------------------
% Theorem environments.
% ---------------------------------------------------------------------------
\theoremstyle{plain}
\newtheorem{theorem}{Theorem}[section]
\newtheorem{lemma}[theorem]{Lemma}
\newtheorem{corollary}[theorem]{Corollary}
\newtheorem{proposition}[theorem]{Proposition}
\theoremstyle{definition}
\newtheorem{definition}[theorem]{Definition}
% Written-track problems are numbered B1, B2, ... on both tracks (P<n> is a
% predict item). Label each one \label{prob:B<n>} and its answer \label{ans:B<n>}.
\newtheorem{problem}{Problem}
\renewcommand{\theproblem}{B\arabic{problem}}
\theoremstyle{remark}
\newtheorem*{remark}{Remark}
\newenvironment{answer}
  {\par\medskip\noindent\textbf{Answer.}\begin{quotation}\noindent}
  {\end{quotation}\medskip}

% ---------------------------------------------------------------------------
% Grader margin notes. Rendered form of the `% GRADER[...][...][...]:` comment
% convention (see template_answers.tex). Use inline where a visible note helps:
%   \gradernote{B2/induction-soundness/2: the invariant is not preserved by the last assignment}
% ---------------------------------------------------------------------------
\definecolor{gradercolor}{RGB}{160, 30, 30}
\newcommand{\gradernote}[1]{\textcolor{gradercolor}{\small\sffamily [GRADER: #1]}}

% ---------------------------------------------------------------------------
% Code listings: Dafny, Verus (Rust inside verus!{...}), Python.
% ---------------------------------------------------------------------------
\definecolor{codekeyword}{RGB}{0, 60, 130}
\definecolor{codecomment}{RGB}{90, 110, 90}
\definecolor{codestring}{RGB}{130, 60, 10}
\definecolor{codehole}{RGB}{180, 0, 0}

\lstdefinelanguage{dafny}{
  sensitive=true,
  morekeywords={
    abstract, allocated, as, assert, assume, bool, break, by, calc, case, char,
    class, codatatype, const, constructor, continue, datatype, decreases, else,
    ensures, exists, expect, export, extends, false, forall, fresh, function,
    ghost, if, import, in, include, int, invariant, is, iterator, label, lemma,
    map, match, method, modifies, modify, module, multiset, nat, new, newtype,
    null, old, opaque, opened, predicate, print, provides, reads, real, refines,
    requires, return, returns, reveal, reveals, seq, set, static, string, then,
    this, trait, true, twostate, type, unchanged, var, while, witness},
  emph={assume},
  emphstyle=\color{codehole}\bfseries,
  comment=[l]{//},
  morecomment=[s]{/*}{*/},
  morestring=[b]",
}

\lstdefinelanguage{verus}{
  sensitive=true,
  morekeywords={
    as, break, const, continue, crate, dyn, else, enum, false, fn, for, if,
    impl, in, let, loop, match, mod, move, mut, pub, ref, return, self, Self,
    static, struct, super, trait, true, type, unsafe, use, where, while,
    spec, proof, exec, open, closed, tracked, ghost, requires, ensures,
    recommends, invariant, decreases, assert, forall, exists, choose, by,
    broadcast, uninterp, axiom, returns, int, nat, Seq, Set, Map, Ghost, Tracked},
  emph={assume, admit},
  emphstyle=\color{codehole}\bfseries,
  comment=[l]{//},
  morecomment=[s]{/*}{*/},
  morestring=[b]",
}

\lstset{
  basicstyle=\ttfamily\small,
  keywordstyle=\color{codekeyword}\bfseries,
  commentstyle=\color{codecomment}\itshape,
  stringstyle=\color{codestring},
  showstringspaces=false,
  columns=fullflexible,
  keepspaces=true,
  breaklines=true,
  frame=leftline,
  framerule=0.6pt,
  xleftmargin=1em,
  aboveskip=0.8em,
  belowskip=0.8em,
  literate=
    {→}{{$\to$}}1        {←}{{$\leftarrow$}}1  {⇒}{{$\Rightarrow$}}1
    {⇔}{{$\Leftrightarrow$}}1                  {↦}{{$\mapsto$}}1
    {∀}{{$\forall$}}1    {∃}{{$\exists$}}1     {λ}{{$\lambda$}}1
    {∧}{{$\wedge$}}1     {∨}{{$\vee$}}1        {¬}{{$\neg$}}1
    {≠}{{$\neq$}}1       {≤}{{$\leq$}}1        {≥}{{$\geq$}}1
    {∈}{{$\in$}}1        {∉}{{$\notin$}}1      {⊆}{{$\subseteq$}}1
    {×}{{$\times$}}1     {·}{{$\cdot$}}1       {∅}{{$\varnothing$}}1
    {ℕ}{{$\mathbb{N}$}}1 {ℤ}{{$\mathbb{Z}$}}1
    {⊢}{{$\vdash$}}1     {⟦}{{$\llbracket$}}1  {⟧}{{$\rrbracket$}}1,
}

% Code environments and inline code.
\lstnewenvironment{dafnycode}{\lstset{language=dafny}}{}
\lstnewenvironment{veruscode}{\lstset{language=verus}}{}
\lstnewenvironment{pycode}{\lstset{language=Python}}{}
\newcommand{\dfy}[1]{\lstinline[language=dafny]{#1}}
\newcommand{\vrs}[1]{\lstinline[language=verus]{#1}}
\newcommand{\py}[1]{\lstinline[language=Python]{#1}}
          % the Makefile always runs latexmk from paper/
%   \begin{document}
%   ...
%   \end{document}
%
% Engine: pdflatex. Dafny and Verus sources are mostly ASCII; the `literate`
% tables below map the common Unicode symbols to math glyphs inside listings.
% For other Unicode, extend a table, or build that file with lualatex
% (`latexmk -lualatex <file>`) after replacing the inputenc/fontenc lines with
% \usepackage{fontspec} and a monospace font with good coverage.
%
% Contents: program-logic macros (Hoare triples, wp, VCs, substitution,
% semantics brackets, typing judgments), inference rules (mathpartir),
% theorem and problem environments, the answer environment, the grader
% margin-note macro, and listings languages `dafny` and `verus` (Python uses
% listings' built-in `Python`), with inline forms \dfy{...}, \vrs{...}, \py{...}.
\usepackage[utf8]{inputenc}
\usepackage[T1]{fontenc}

\usepackage{amsmath}
\usepackage{amssymb}
\usepackage{amsthm}
\usepackage{stmaryrd}         % \llbracket \rrbracket for semantics brackets
\usepackage{mathpartir}       % inference rules: \inferrule, the mathpar environment
\usepackage{listings}
\usepackage{xcolor}
\usepackage[margin=1.1in]{geometry}

% ---------------------------------------------------------------------------
% Program logic.
%   \hoare{P}{c}{Q}       {P} c {Q}           (partial correctness)
%   \thoare{P}{c}{Q}      [P] c [Q]           (total correctness)
%   \provable{P}{c}{Q}    |- {P} c {Q}        (derivable in the proof system)
%   \WP{c}{Q}             wp(c, Q)            (\wp is LaTeX's Weierstrass p, so
%   \WLP{c}{Q}            wlp(c, Q)            the macro names are upper case)
%   \VC{c}{Q}             VC(c, Q)
%   \subst{Q}{x}{e}       Q[e/x]              (e replaces x in Q)
%   \sem{e}               [[e]]               (denotation)
%   \assign, \seqc        :=  and the sequencing semicolon
%   \impl                 ==> as a math arrow (LaTeX's own \iff gives <==>)
% ---------------------------------------------------------------------------
\newcommand{\hoare}[3]{\{\,#1\,\}\;#2\;\{\,#3\,\}}
\newcommand{\thoare}[3]{[\,#1\,]\;#2\;[\,#3\,]}
\newcommand{\provable}[3]{\vdash \hoare{#1}{#2}{#3}}
\newcommand{\WP}[2]{\mathrm{wp}(#1,\,#2)}
\newcommand{\WLP}[2]{\mathrm{wlp}(#1,\,#2)}
\newcommand{\VC}[2]{\mathrm{VC}(#1,\,#2)}
\newcommand{\subst}[3]{#1[#3/#2]}
\newcommand{\sem}[1]{\llbracket #1 \rrbracket}
\newcommand{\assign}{\mathrel{:=}}
\newcommand{\seqc}{\mathbin{;}\,}
\newcommand{\impl}{\Longrightarrow}

% ---------------------------------------------------------------------------
% Judgments.  \judgment{\Gamma}{e}{\tau}  ==>  Gamma |- e : tau
% ---------------------------------------------------------------------------
\newcommand{\judgment}[3]{#1 \vdash #2 : #3}
\newcommand{\judg}[3]{\judgment{#1}{#2}{#3}}          % short alias
\newcommand{\emptyjudgment}[2]{\cdot \vdash #1 : #2}  % empty context
\newcommand{\ctx}{\Gamma}
\newcommand{\turnstile}{\vdash}

% ---------------------------------------------------------------------------
% Theorem environments.
% ---------------------------------------------------------------------------
\theoremstyle{plain}
\newtheorem{theorem}{Theorem}[section]
\newtheorem{lemma}[theorem]{Lemma}
\newtheorem{corollary}[theorem]{Corollary}
\newtheorem{proposition}[theorem]{Proposition}
\theoremstyle{definition}
\newtheorem{definition}[theorem]{Definition}
% Written-track problems are numbered B1, B2, ... on both tracks (P<n> is a
% predict item). Label each one \label{prob:B<n>} and its answer \label{ans:B<n>}.
\newtheorem{problem}{Problem}
\renewcommand{\theproblem}{B\arabic{problem}}
\theoremstyle{remark}
\newtheorem*{remark}{Remark}
\newenvironment{answer}
  {\par\medskip\noindent\textbf{Answer.}\begin{quotation}\noindent}
  {\end{quotation}\medskip}

% ---------------------------------------------------------------------------
% Grader margin notes. Rendered form of the `% GRADER[...][...][...]:` comment
% convention (see template_answers.tex). Use inline where a visible note helps:
%   \gradernote{B2/induction-soundness/2: the invariant is not preserved by the last assignment}
% ---------------------------------------------------------------------------
\definecolor{gradercolor}{RGB}{160, 30, 30}
\newcommand{\gradernote}[1]{\textcolor{gradercolor}{\small\sffamily [GRADER: #1]}}

% ---------------------------------------------------------------------------
% Code listings: Dafny, Verus (Rust inside verus!{...}), Python.
% ---------------------------------------------------------------------------
\definecolor{codekeyword}{RGB}{0, 60, 130}
\definecolor{codecomment}{RGB}{90, 110, 90}
\definecolor{codestring}{RGB}{130, 60, 10}
\definecolor{codehole}{RGB}{180, 0, 0}

\lstdefinelanguage{dafny}{
  sensitive=true,
  morekeywords={
    abstract, allocated, as, assert, assume, bool, break, by, calc, case, char,
    class, codatatype, const, constructor, continue, datatype, decreases, else,
    ensures, exists, expect, export, extends, false, forall, fresh, function,
    ghost, if, import, in, include, int, invariant, is, iterator, label, lemma,
    map, match, method, modifies, modify, module, multiset, nat, new, newtype,
    null, old, opaque, opened, predicate, print, provides, reads, real, refines,
    requires, return, returns, reveal, reveals, seq, set, static, string, then,
    this, trait, true, twostate, type, unchanged, var, while, witness},
  emph={assume},
  emphstyle=\color{codehole}\bfseries,
  comment=[l]{//},
  morecomment=[s]{/*}{*/},
  morestring=[b]",
}

\lstdefinelanguage{verus}{
  sensitive=true,
  morekeywords={
    as, break, const, continue, crate, dyn, else, enum, false, fn, for, if,
    impl, in, let, loop, match, mod, move, mut, pub, ref, return, self, Self,
    static, struct, super, trait, true, type, unsafe, use, where, while,
    spec, proof, exec, open, closed, tracked, ghost, requires, ensures,
    recommends, invariant, decreases, assert, forall, exists, choose, by,
    broadcast, uninterp, axiom, returns, int, nat, Seq, Set, Map, Ghost, Tracked},
  emph={assume, admit},
  emphstyle=\color{codehole}\bfseries,
  comment=[l]{//},
  morecomment=[s]{/*}{*/},
  morestring=[b]",
}

\lstset{
  basicstyle=\ttfamily\small,
  keywordstyle=\color{codekeyword}\bfseries,
  commentstyle=\color{codecomment}\itshape,
  stringstyle=\color{codestring},
  showstringspaces=false,
  columns=fullflexible,
  keepspaces=true,
  breaklines=true,
  frame=leftline,
  framerule=0.6pt,
  xleftmargin=1em,
  aboveskip=0.8em,
  belowskip=0.8em,
  literate=
    {→}{{$\to$}}1        {←}{{$\leftarrow$}}1  {⇒}{{$\Rightarrow$}}1
    {⇔}{{$\Leftrightarrow$}}1                  {↦}{{$\mapsto$}}1
    {∀}{{$\forall$}}1    {∃}{{$\exists$}}1     {λ}{{$\lambda$}}1
    {∧}{{$\wedge$}}1     {∨}{{$\vee$}}1        {¬}{{$\neg$}}1
    {≠}{{$\neq$}}1       {≤}{{$\leq$}}1        {≥}{{$\geq$}}1
    {∈}{{$\in$}}1        {∉}{{$\notin$}}1      {⊆}{{$\subseteq$}}1
    {×}{{$\times$}}1     {·}{{$\cdot$}}1       {∅}{{$\varnothing$}}1
    {ℕ}{{$\mathbb{N}$}}1 {ℤ}{{$\mathbb{Z}$}}1
    {⊢}{{$\vdash$}}1     {⟦}{{$\llbracket$}}1  {⟧}{{$\rrbracket$}}1,
}

% Code environments and inline code.
\lstnewenvironment{dafnycode}{\lstset{language=dafny}}{}
\lstnewenvironment{veruscode}{\lstset{language=verus}}{}
\lstnewenvironment{pycode}{\lstset{language=Python}}{}
\newcommand{\dfy}[1]{\lstinline[language=dafny]{#1}}
\newcommand{\vrs}[1]{\lstinline[language=verus]{#1}}
\newcommand{\py}[1]{\lstinline[language=Python]{#1}}


\title{Q12 --- Practice I-C + Lab 1 \\ \large Written answers}
\author{learner}
\date{}

\begin{document}
\maketitle

% PREREG-BEGIN
\section*{Problem B1 (pre-registration)}\label{ans:B1}
\begin{answer}
\textbf{Claim} (the three hypotheses, A2 vs A1, A3 vs A1 and the language
interaction, each one- or two-sided): \dots

\textbf{Metric and estimator} (per task and repeat: solved within $B$ calls; the
arm's rate; the paired difference per comparison): \dots

\textbf{Arms and policies} (the three Qwen arms and the Claude arm's pick, each
with its policy file and prompt hash; the Claude subset's seed and its
stratification): \dots

\textbf{Tasks and strata} (the sets' commits, the validation files' sha256, the
draw seed and the members it draws from: each pool less the tasks another lab's
\texttt{[[draw]]} record took, excluding Lab 2's HumanEval problems across both languages
(the five Verus tasks and four matching problems from the Dafny stratum; D39); the pilot rule; the local stratum's composition in each case): \dots

\textbf{Budget arithmetic} (calls = tasks $\times$ attempts $\times$ rounds
$\times$ arms $\times$ repeats, per arm, plus the sentinel calls; the
throughput assumption and its source; against the registry's budget): \dots

\textbf{Generation settings} (the values the Ollama gate pins for every Qwen
call; the seed list per arm and for the pilot, pairwise disjoint; for Claude
the pinned argv): \dots

\textbf{Cost per arm} (tokens in and out, GPU-seconds and wall clock per arm,
against the ceiling; the truncation rate per arm as a sensitivity line): \dots

\textbf{MDE} (computed with \texttt{fvh.stats.mde\_two\_proportions}: $n$,
baseline, pairing, Holm-adjusted $\alpha$, the call and its value): \dots

\textbf{Seeds and repeats} (Qwen: the per-request seeds; Claude: repeated
runs): \dots

\textbf{Cheat surface and gate} (the instruments every claimed success passes,
the X classes, \texttt{UNKNOWN} handling): \dots

\textbf{Exclusions} (\texttt{infra-failed} re-issue; \texttt{truncated:<cause>}
as failures, with the sensitivity line): \dots

\textbf{Stopping rule}: \dots

\textbf{Prediction intervals} (80\%, for each arm's rate and each difference,
scored after the run): \dots

\textbf{Contamination} (the local stratum's authorship dates; the public sets
flagged as possibly seen): \dots

\textbf{Disclosure} (a small-scale replication with model substitution; Claude as
grader and subject; one machine, one solver pin): \dots

\textbf{Batch plan} (the plan file \texttt{python -m fvh run Q12 -{}-plan}
reads: its path, and its sha256 on a line of its own as \texttt{plan-sha256:}
followed by the 64 hex digits \texttt{sha256sum} prints for the file;
\texttt{fvh run} opens a lab batch only on the plan this line names;
separately executed pilot or subset plans need their own binding and a new
gradebook row after re-grading the changed pre-registration, with the
registered method held fixed): \dots
\end{answer}
% PREREG-END

\section*{Problem B2 (lab report)}\label{ans:B2}
\begin{answer}
\textbf{Reproduce:} \verb+<one command over research/runs/Q12/ that regenerates every table below>+

\textbf{Results} (the paired tables and comparisons, with intervals): \dots

\textbf{Deviations from the pre-registration} (each one, and why): \dots

\textbf{Transcript and cheat audit} (the gate's verdict on every claimed
success; \texttt{UNKNOWN} classes): \dots

\textbf{Prediction intervals, scored}: \dots

\textbf{The hand part beside the arms} (the shared local tasks; descriptive
only): \dots

\textbf{Variance and intervals}: \dots

\textbf{Contamination}: \dots

\textbf{Disclosure}: \dots
\end{answer}

% ---------------------------------------------------------------------------
\section*{Problem B3 (paper card)}\label{ans:B3}
\begin{answer}
\textbf{Source:} cheatsheet entry 05.10, lines 2148--2171 (the pinned range).

\textbf{Task shape:} \dots

\textbf{Oracle:} \dots

\textbf{What is trusted:} \dots

\textbf{Cheat surface:} \dots

\textbf{Metric:} \dots

\textbf{Budget:} \dots

\textbf{Variance:} \dots

\textbf{Artifact:} \dots

\textbf{Closest arm, and what its budget counts:} \dots
\end{answer}

% ---------------------------------------------------------------------------
\section*{Problem B4 (paper card)}\label{ans:B4}
\begin{answer}
\textbf{Source:} cheatsheet entry 06.1, lines 2312--2338 (the pinned range).

\textbf{Task shape:} \dots

\textbf{Oracle:} \dots

\textbf{What is trusted:} \dots

\textbf{Cheat surface:} \dots

\textbf{Metric:} \dots

\textbf{Budget:} \dots

\textbf{Variance:} \dots

\textbf{Artifact:} \dots

\textbf{Closest arm, and what its budget counts:} \dots
\end{answer}

% ---------------------------------------------------------------------------
\section*{Problem B5 (paper card)}\label{ans:B5}
\begin{answer}
\textbf{Source:} cheatsheet entry 04.8, lines 1681--1706 (the pinned range).

\textbf{Task shape:} \dots

\textbf{Oracle:} \dots

\textbf{What is trusted:} \dots

\textbf{Cheat surface:} \dots

\textbf{Metric:} \dots

\textbf{Budget:} \dots

\textbf{Variance:} \dots

\textbf{Artifact:} \dots

\textbf{Closest arm, and what its budget counts:} \dots
\end{answer}

\end{document}
